\itemhead{习题3.2 · 第8题（教材p.86）}{XJ-B1-C03-HW-P086-Y0520}
\label{XJ-B1-C03-HW-P086-Y0520}
{\footnotesize 湘教版教材p.86 / PDF第93页 / 习题3.2 / 第8题}
\field{原题}{设函数 $f(x)$ 的定义域为 $(-4,5)$，如果 $f(x)$ 在 $(-4,0)$ 上是减函数，在 $(0,5)$ 上也是减函数，能不能断定它在 $(-4,5)$ 上是减函数？如果 $f(x)$ 在 $(-4,0]$ 上是增函数，在 $[0,5)$ 上也是增函数，能不能断定它在 $(-4,5)$ 上是增函数？}
\field{参考答案}{第一问不能。例如 $f(x)=-x$（$x<0$）、$f(x)=100-x$（$x>0$），任取 $f(0)$，两侧递减但跨0不递减。第二问能：跨0时由 $f(x_1)<f(0)<f(x_2)$ 推出，且两段包含公共点0。}
{\small \field{关联知识}{F05 单调性定义、证明与参数（教材帮p.97、98、100 / PDF106、107、109）}
\field{方法与易错点}{用定义比较任意两值，或依据已知图象性质；含参时使指定区间被相应单调区间覆盖。 易错：用有限采样代替证明；把分离半轴合并；漏跨零点情形。}
\field{答案依据}{答案册相关题组：PDF第23、24页（册内页22、23）；按当前题库原题逐项重算，不按旧版题号机械对套。}}
